\begin{exercisechunk}
\item \textbf{A robust separation from alternating tent-map labels.}
\label[exercise]{ex:l14-tent-robust-separation}
\exerciseprofile{3}{2}{3}{\exproof\quad\excalculation\quad\exinterpretation}{%
Extend the exact folding separation to a quantitative classification-error
lower bound.}
The lecture proves that an exact shallow representation of an iterated tent
map needs exponentially many hinges. This exercise shows that a shallow
network also needs exponentially many pieces to obtain small error on a
specified finite classification problem.

Fix \(L\ge1\), put \(N=2^L\), and define
\[
 x_j=\frac{j+1/4}{N},
 \qquad
 y_j=\I\bigl(\tau^{\circ L}(x_j)\ge1/2\bigr),
 \qquad j=0,\ldots,N-1.
\]
For a scalar function \(f\), define
\[
 \widehat y_j(f)=\I(f(x_j)\ge1/2),
 \qquad
 \operatorname{err}_L(f)
 =\frac1N\sum_{j=0}^{N-1}\I\bigl(\widehat y_j(f)\ne y_j\bigr).
\]
\begin{enumerate}[(a),beginpenalty=10000]
\item Use the slopes and endpoint values of \(\tau^{\circ L}\) to prove
\[
 \tau^{\circ L}(x_j)
 =\begin{cases}1/4,&j\text{ even},\\[1mm]3/4,&j\text{ odd},\end{cases}
 \qquad
 y_j=j\bmod2.
\]
\item Let \(\widehat y_0,\ldots,\widehat y_{N-1}\) be any binary sequence
with \(R\) maximal constant runs. Prove that its number of disagreements with
the alternating sequence \(0,1,0,1,\ldots\) is at least \((N-R)/2\).
\item Suppose \(f\) is continuous and affine on at most \(K\) intervals.
Prove that \(\widehat y_0(f),\ldots,\widehat y_{N-1}(f)\) has at most \(2K\)
constant runs. Deduce
\[
 \operatorname{err}_L(f)
 \ge \frac12-\frac{K}{2^L}.
\]
\item A one-hidden-layer scalar ReLU network with \(m\) hidden units is affine
on at most \(m+1\) intervals. Prove that
\[
 \operatorname{err}_L(f)<\frac14
 \quad\Longrightarrow\quad
 m+1>2^{L-2}.
\]
Compare this necessary shallow width with the \(O(L)\)-unit deep
representation of \(\tau^{\circ L}\).
\end{enumerate}
\begin{hints}
For part~(b), count the errors inside each constant run separately. For
part~(c), the superlevel set of an affine function on one interval is empty,
the full interval, or an interval touching one endpoint.
\end{hints}
\end{exercisechunk}
